% =======================================================
% 2. DOCUMENTAZIONE
% =======================================================
\section{Documentazione}

\subsection{Template e convenzioni}
[-]

\subsection{Struttura documenti}
[-]

\subsection{Ciclo di vita documentale}
\begin{enumerate}
	\item \textbf{Redazione}: stesura del documento da parte di un membro;
	\item \textbf{Verifica}: controllo incrociato da parte di un verificatore (diverso dal redattore);
	\item \textbf{Approvazione}: rilascio ufficiale autorizzato dal Responsabile.
\end{enumerate}

\subsection{Strumenti di documentazione}
Il gruppo utilizza LaTeX (tramite TeXstudio/Overleaf) per la stesura dei documenti, strutturati in cartelle logiche.

\subsection{Gestione della Configurazione}
[-]

\subsection{Versionamento}
Viene adottato un sistema basato sul Semantic Versioning a tre cifre: \textbf{vX.Y.Z}
\begin{itemize}
	\item \textbf{X (Major)}: rilasci ufficiali per le baseline o cambiamenti strutturali profondi;
	\item \textbf{Y (Minor)}: aggiunta di nuovi capitoli o funzionalità corpose;
	\item \textbf{Z (Patch)}: correzioni ortografiche o piccoli bug fix.
\end{itemize}

\subsection{Nomenclatura file e branch}
[-]

\subsection{Repository e struttura}
[-]

\subsection{Gestione della Qualità}
[-]

\subsection{Obiettivi di qualità}
[-]

\subsection{Metriche}
[-]

\subsection{Verifica}
[-]

\subsection{Analisi statica}
[-]

\subsection{Analisi dinamica}
[-]

\subsection{Walkthrough e inspection}
\begin{itemize}
	\item \textbf{Walkthrough}: [-]
	\item \textbf{Inspection}: [-]
\end{itemize}

\subsection{Validazione}
[-]